\documentclass{article}
\usepackage{microtype}
\usepackage{graphicx}
\usepackage{booktabs}
\usepackage{amsmath}
\usepackage{hyperref}
\usepackage{tikz}
\usetikzlibrary{arrows.meta,positioning,fit}
% Anonymous submission. For the camera-ready version use: \usepackage[accepted]{icml2026}
\usepackage{icml2026}
\icmltitlerunning{What Does a Matching Network Statistic Establish?}

\renewcommand{\dbltopfraction}{.92}\renewcommand{\dblfloatpagefraction}{.9}\renewcommand{\textfraction}{.05}
\begin{document}
\twocolumn[
\icmltitle{What Does a Matching Network Statistic Establish?\\
Procedure, Density, and Sampling in LLM-Generated Friendship Networks}
\begin{icmlauthorlist}
\icmlauthor{Anonymous Authors}{anon}
\end{icmlauthorlist}
\icmlaffiliation{anon}{Anonymous Institution}
\icmlcorrespondingauthor{Anonymous Author}{anonymous@example.invalid}
\icmlkeywords{network generation, social simulation, model evaluation, homophily}
\vskip 0.3in
]
\printAffiliationsAndNotice{}

\begin{abstract}
When a language model generates a friendship network that matches an observed one on some statistic, what has been shown? We ask four questions of one design: whether prompt framing changes the networks, whether the generation procedure changes which attribute organises ties, whether a structural match reflects more than edge count, and whether the verdict depends on how the observed network was sampled. Four model configurations generate networks over a fixed roster of 50 fictional adults under four procedures and seven country-framing or instruction-language settings; 890 of 896 planned attempts yield valid graphs and all analyses reuse them. Framing barely matters: one of 24 planned contrasts passes familywise correction, and ties generated under a different setting overlap as much as ties from repeated runs of the same setting (97--101\% of the within-setting overlap). Procedure matters: person-by-person nomination sorts ties most strongly by political orientation and whole-network generation by religion, in every model--setting cell under a permutation reference that conditions on the other attributes, and in a joint dyadic model with eight of eight schedule blocks agreeing in sign. Against eleven observed records, whole-network generation is closest on global efficiency, but a random-edge generator given only the target density closes 87\% of the gap while remaining the worst candidate on clustering. Replacing uniform with snowball samples of the same records moves whole-network generation from closest on efficiency in nine records to closest in none. A matching statistic is evidence about a generator only after procedure, density and sampling have been ruled out as explanations.
\end{abstract}

\section{Introduction}
Two friendship graphs can have the same number of edges and quite different social organisation. The same ties can be concentrated in a few tightly connected groups or spread across groups, shortening the paths between them. A generator that matches an observed graph on path length may therefore have captured something about where friendships fall, or it may simply have drawn approximately the right number of edges. There is a second ambiguity on the observed side: when only 50 people are sampled from a larger network, a friendship disappears whenever either endpoint falls outside the sample. Which of these properties is a comparison rewarding?

Language models make these questions testable because the same fictional people can be presented to a model under different rules for choosing ties. \citet{chang} found that person-level and whole-network prompting produce different network structures and that generated ties overstate political homophily. \citet{kilaru} compared four generation procedures across cultural and linguistic settings in a 192-network directed study and also reported political sorting under three of them. We do not introduce those procedures, and we do not pool those earlier graphs with ours. We ask instead whether conclusions about mixing persist across scoring rules and matched repetitions, and whether an apparent structural advantage remains once a simple generator is given the target density.

We hold 50 biographies fixed (Figure~\ref{fpipe}). Four model configurations either generate a complete network at once (\emph{Global}), nominate friends one person at a time (\emph{Local}), nominate in order with information about earlier decisions (\emph{Sequential}), or revise initial nominations over three rounds (\emph{Iterative}). We vary country framing with English instructions and instruction language with US framing. Eight matched schedule blocks let us separate the direction of a prompt contrast from its uncertainty. Six failed attempts are kept as failures and twelve valid empty graphs as empty; six graphs produced in a later supplement never replace missing observations and enter no analysis.

\begin{figure*}[t]\centering
\begin{tikzpicture}[font=\scriptsize,>={Stealth[length=1.6mm]},
 box/.style={draw=black!65,rounded corners=1.5pt,align=center,inner sep=3pt,minimum height=9.5mm,text width=#1},
 box/.default=27mm,
 gen/.style={box=#1,fill=blue!6},obs/.style={box=#1,fill=orange!10},ana/.style={box=#1,fill=black!5},
 diag/.style={box=#1,fill=white,dashed}]
\node[gen=26mm](roster){\textbf{Fixed roster}\\50 fictional adults;\\politics, religion,\\gender};
\node[gen=38mm,right=5mm of roster](design){\textbf{Generation design}\\4 configurations $\times$ 4 procedures\\$\times$ 7 settings $\times$ 8 blocks $=$ 896 attempts};
\node[gen=33mm,right=5mm of design](valid){\textbf{Validation}\\890 valid graphs (12 empty)\\6 failures kept as failures};
\node[ana=30mm,right=7mm of valid,yshift=9mm](contrast){\textbf{Planned contrasts}\\24 paired contrasts,\\Holm and Bonferroni (\S\ref{sec:contrasts})};
\node[ana=30mm,below=2.5mm of contrast](mix){\textbf{Mixing}\\3 scores, 28 cells;\\stratified permutation (\S\ref{sec:mixing}, \ref{sec:robust})};
\node[ana=30mm,below=2.5mm of mix](ext){\textbf{External comparison}\\$W_1$ on six statistics,\\50-node draws (\S\ref{sec:density})};
\node[obs=38mm,below=6mm of design](records){\textbf{Observed records}\\10 Facebook, 1 school (test);\\8 development records; Copenhagen};
\node[obs=33mm,below=6mm of valid](base){\textbf{Fixed baselines}\\ER, BA, WS fitted on\\development records only};
\node[diag=24mm,below=6mm of roster](d1){\textbf{Diagnostics}\\target-density ER;\\uniform vs.\ snowball};
\draw[->](roster)--(design);\draw[->](design)--(valid);
\draw[->](valid.east)--++(3mm,0)|-(contrast.west);
\draw[->](valid.east)--++(3mm,0)|-(mix.west);
\draw[->](valid.east)--++(3mm,0)|-(ext.west);
\draw[->](records)--(base);
\draw[->](base.east)--++(3mm,0)|-([yshift=-2mm]ext.west);
\draw[->,dashed](d1)--(records);
\end{tikzpicture}
\caption{Study pipeline. Blue: generation, with every planned attempt retained in the denominator. Orange: observed data; baselines are fixed on development records before any test record is seen. Grey: the three analyses. Dashed: post hoc diagnostics that alter either the information given to a baseline or the sampling of observed records (Sections~\ref{sec:density}--\ref{sec:robust}); they use no new generation.}\label{fpipe}\end{figure*}

\paragraph{Research questions and answers.}
\begin{itemize}
\setlength{\itemsep}{1pt}
\item \textbf{RQ1. Do country framing and instruction language change the generated networks?} \emph{Hardly.} Of 24 planned contrasts, 23 were testable and one passed correction: Iterative US-Hindi versus US-English clustering, significant in only three of eight leave-one-block-out refits. At the level of individual ties, runs under a different setting share as many ties with the reference as repeated reference runs share with each other (Sections~\ref{sec:contrasts} and~\ref{sec:robust}).
\item \textbf{RQ2. Does the generation procedure change which attribute organises ties?} \emph{Yes, and it reverses the order.} Politics is the most sorted attribute under person-level procedures and religion under Global, under a permutation reference that holds the other two attributes fixed (28 of 28 cells each) and in a joint dyadic model (Sections~\ref{sec:mixing} and~\ref{sec:robust}).
\item \textbf{RQ3. Is a close structural match evidence about the generator or about its edge count?} \emph{Mostly edge count for reachability, not for closure.} Giving a random-edge generator each record's density closes 87\% of its efficiency gap to Global over eleven records, yet it remains the worst candidate on clustering and transitivity (Sections~\ref{sec:density} and~\ref{sec:robust}).
\item \textbf{RQ4. Does the verdict depend on how the observed network was sampled?} \emph{Yes.} Uniform 50-node draws from a Copenhagen friendship network keep density and lose clustering and efficiency; on the eleven test records, snowball draws triple density and reverse the efficiency ranking of the candidates (Sections~\ref{sec:cns} and~\ref{sec:robust}).
\end{itemize}
None of these results establishes differences or equivalence between human populations, and none validates prediction of individual friendships.

\section{Related work}
\citet{chang} compared prompted friendship graphs with observed networks and reported excess political homophily. We extend that evaluation in two directions: whether the same attribute dominates under different nomination rules and chance references, and whether a structural match survives a baseline informed only by tie density. The directed study of \citet{kilaru} supplies procedural context, not additional observations; its Global comparison included age, which our three-attribute ranking does not test.

Work on agent interaction examines how network structure and opinions evolve together \citep{piao,papachristou}. We condition on fixed biographies: there is no opinion trajectory, so assortativity here does not measure polarisation over time. Multilingual political assessments distinguish language from nationality framing \citep{helwe}; we apply those two changes to friendship instructions, not survey responses. Simulating a human sample requires a substantive account of correspondence with the sampled population \citep{argyle}.

Discrepancy metrics for graph generators are sensitive to how the comparison is constructed \citep{obray}, and embedding-based alternatives assess distributional fidelity and diversity \citep{thompson}. We deliberately estimate marginal distances for named statistics, not agreement on the joint distribution of graph properties. Sampling preserves some network properties and distorts others \citep{sampling}; we therefore treat the sampling scheme as part of the estimand.

\begin{figure*}[t]\centering
\includegraphics[width=.86\linewidth,height=.44\textheight,keepaspectratio]{01_primary_contrasts.pdf}
\caption{All 24 planned contrasts (setting minus US-English within the same seed block). Circles are means over complete blocks, lines are Bonferroni simultaneous 95\% intervals for the family of 24, and small dots are the individual paired blocks. The contrast with five complete blocks was not tested and has no interval.}\label{f01}\end{figure*}

\section{Design and estimands}
\paragraph{Roster and schedule.} The roster has 50 fictional adults described in English, with political orientation (17 left-leaning, 17 right-leaning, 16 centrist), religion (five groups of 10) and gender (20 women, 20 men, 10 nonbinary); age enters only the dyadic model of Section~\ref{sec:robust}. Four model configurations (release identifiers GPT-4.1, GPT-5.6-Luna, GPT-6-Luna and GPT-6-Sol) are each run under the four procedures. The reference setting is US framing with English instructions. Three settings change the framing to India, Japan or Brazil in English; three change the instruction language to Hindi, Japanese or Portuguese under US framing. This seven-arm design isolates each intervention relative to one reference; it cannot estimate country-by-language interactions. Eight schedule repetitions give $4\times4\times7\times8=896$ planned attempts. The unit of analysis is the resulting simple undirected graph on the roster, not a model response, a conversational turn or an edge.

\begin{table}[t]\centering
\caption{Analysis populations. \emph{Valid} includes graphs with no ties; \emph{nonempty} is the denominator for mixing scores. Supplementary attempts never enter any estimate.}\label{tab:coverage}
\small
\begin{tabular}{lrrr}\toprule Procedure & Planned & Valid & Nonempty\\\midrule
Global &224&219&207\\Local &224&224&224\\Sequential &224&224&224\\Iterative &224&223&223\\\midrule Total&896&890&878\\\bottomrule\end{tabular}\end{table}

\paragraph{Planned comparisons.} For schedule block $b$, let $d_b$ be the mean across the four configurations of their paired setting-minus-reference differences. The primary contrast is the mean of $d_b$ over blocks in which every required pair is defined. The endpoint is edge density for Global and mean local clustering for the person-level procedures. We require at least six complete blocks, use two-sided paired-$t$ inference on block differences, report Bonferroni simultaneous 95\% intervals, and apply Holm's step-down adjustment \citep{holm} across the full planned family of 24; an untestable contrast stays in the family. The intervals and the adjusted $p$-values come from different procedures and coincide only for the smallest $p$-value. Model-specific profiles are descriptive and are not a second family of tests.

\paragraph{Mixing scores.} For attribute $a$, let $s_a$ be the fraction of edges whose endpoints share a category. With $n_g$ people in category $g$ and $n=50$, the roster-pair reference is $q_a=\sum_g n_g(n_g-1)/[n(n-1)]$ (0.320 for politics, 0.184 for religion, 0.347 for gender) and the \emph{excess share} is $s_a-q_a$. Let $p_g=n_g/n$, let $a_g$ be the share of edge ends attached to category $g$, and let $h_g$ be the share of those ends whose other end is also in $g$. The \emph{Coleman-type index} is $\sum_g p_g (h_g-p_g)/(1-p_g)$, a size-weighted version of the normalisation of \citet{coleman}; the classical piecewise index gives the same rank counts on these graphs. \emph{Newman's assortativity} is $(s_a-\sum_g a_g^2)/(1-\sum_g a_g^2)$ \citep{newman}, which centres on edge-end shares, so the references differ when categories differ in degree. We rank attributes within each model--setting cell using repetitions on which all scores are defined. For contrasts between procedures we first pair on model, setting and schedule, average within the cell, and then weight the 28 cells equally. No score is imputed to a graph without edges.

\begin{figure*}[t]\centering
\includegraphics[width=\linewidth,height=.40\textheight,keepaspectratio]{14_demographic_mixing.pdf}
\caption{Demographic mixing in the 878 nonempty graphs. Left: share of ties within a group minus the roster-chance share, by procedure and attribute. Right: politics score minus the larger of the religion and gender scores under three scores, with the number of cells in which politics ranks first. Each tick is one of 28 model--setting cells; tick spread is not a confidence interval.}\label{f14}\end{figure*}

\paragraph{External comparison.} Baselines were fixed on eight archived friendship records from seven communities, averaging the two waves of one community before fitting. The development mean density sets $p=0.189$ for independent random edges (ER; \citealt{gilbert}); matching that density sets $m=5$ for preferential attachment (BA; \citealt{ba}) and $k=10$ for small-world graphs (WS; \citealt{ws}), whose rewiring probability of 0.2 was chosen by a grid search on development clustering. Each baseline has eight simulation banks of 56 graphs. Nothing is refitted on test outcomes.

The transfer panel consists of ten Facebook alter networks with the ego removed \citep{mcauley} and one school friendship survey \citep{school}, in which a nomination in either direction between roster members is a tie. We take 32 uniform 50-node induced draws per record and keep disconnected and empty draws. For each statistic we compute the one-dimensional Wasserstein distance $W_1=\int|F(x)-G(x)|\,dx$ between the candidate and target distributions; smaller is closer on that statistic only. Within a model--procedure group each instruction setting has equal weight, conditional on a valid graph. We average the four model-specific distances for a procedure and weight records equally when pooling. Ranges over simulation banks or over draws describe variation; they are not confidence intervals.

We evaluate density, mean local clustering, transitivity, global efficiency (mean reciprocal shortest-path length, with disconnected pairs contributing zero; \citealt{latora}), the Gini coefficient of degree, and the fraction of nodes in the largest component. A separate \emph{target-informed diagnostic} gives ER the density of each full observed record. Unlike a transfer candidate it knows a property of the target before generating, so it is never ranked against them.

\begin{figure*}[t]\centering
\includegraphics[width=.9\linewidth,height=.29\textheight,keepaspectratio]{16_rq2_paired_mixing.pdf}
\caption{Paired differences between procedures in excess same-group share, matched on model, setting and schedule before averaging. Each tick is one cell; circles are equal-weight means of the 28 cells, with the number of cells sharing the sign of the mean. Pairs with an empty or failed graph are missing, not replaced.}\label{f16}\end{figure*}

\begin{figure*}[t]\centering
\includegraphics[width=.8\linewidth,height=.29\textheight,keepaspectratio]{15_efficiency_density.pdf}
\caption{Efficiency $W_1$ by observed record, ordered by full-record density. Left: Global, the fixed WS baseline and the density-informed ER diagnostic. Right: Global minus the diagnostic against record density. The diagnostic uses information unavailable to transfer candidates.}\label{f15}\end{figure*}

\section{Results}
\subsection{RQ1: coverage and planned contrasts}\label{sec:contrasts}
Table~\ref{tab:coverage} gives the analysis populations. The six failures fall in three procedure--setting groups of 32 scheduled attempts each: three invalid Global US-Japanese responses, two invalid Global US-Portuguese responses, and one Iterative US-Japanese run that exhausted its corrections. The affected contrasts have five, six and seven complete blocks; the other 21 have eight. For Global US-Japanese the five-block density difference is $+0.0036$. Letting each missing endpoint range over $[0,1]$ bounds the eight-block difference to $[+0.0010,+0.0948]$; this is an identification bound, not a confidence interval. Twelve valid graphs without ties remain in every structural summary for which the endpoint is defined.

Figure~\ref{f01} and Appendix Table~\ref{contrasts} report the whole family. None of the twelve country-framing contrasts passes correction. For Iterative, US-Hindi instructions lower mean clustering by 0.0453. Omitting any one block leaves a negative estimate ($-0.0480$ to $-0.0376$) with raw $p\le0.0036$, but only three of the eight refits pass Holm. A post hoc Shapiro--Wilk test on the eight differences gives $p=0.013$, so the $t$ reference distribution is itself uncertain. An exact sign-flip test cannot arbitrate: with eight blocks its smallest two-sided $p$-value is $2/256$, which exceeds the first Holm threshold $0.05/24$, so its adjusted value of 0.1875 is the floor of the test and not evidence against the effect. All four configuration-specific differences are negative ($-0.022$ to $-0.070$) on the common blocks; they describe one design, not four replications. Failing to reject the other contrasts is not evidence that their effects are negligible.

\begin{figure*}[t]\centering
\includegraphics[width=\linewidth,height=.38\textheight,keepaspectratio]{18_sampling_distortion.pdf}
\caption{Full observed records (open diamonds) against induced draws of 25, 40 and 50 nodes, for twelve records ordered by size. Points are means of 32 overlapping draws and lines their full range; ranges are sampling variation, not confidence intervals.}\label{f18}\end{figure*}

\subsection{RQ2: political mixing ranks first under person-level procedures}\label{sec:mixing}
Politics ranks first by excess same-group share in 85 of 112 cells: all 84 person-level cells and one Global cell (Figure~\ref{f14}). Under both normalised scores it ranks first in 90 cells, the five additional cells all being Global. Removing each schedule in turn leaves every person-level procedure at 28 of 28 for every score. Cell-level consistency does not imply that each individual graph has the same ranking.

The change in rank comes mostly from religion, not politics (Figure~\ref{f16}). Compared with Global, mean excess same-religion share is lower by 0.311, 0.310 and 0.406 for Local, Sequential and Iterative, with all 28 paired cell means in that direction for each comparison. Sequential minus Global changes political excess by only $+0.008$. A switch in the leading attribute therefore need not mean stronger political sorting. Iterative does sort more on politics than Sequential ($+0.129$ over 223 defined pairs, positive in all 28 cells). Appendix Figure~\ref{f17} shows which group pairs carry these changes.

Three features of the design limit interpretation. Religion and politics are associated on this one roster (Cram\'er's $V=0.316$; Section~\ref{sec:robust} conditions on this association), the order of attribute fields in the prompt was fixed, and procedures used very different numbers of requests. The scores describe mixing in generated ties; they do not isolate the causal importance of an attribute or any human preference.

\subsection{RQ3: what a close structural match buys}\label{sec:density}
On global efficiency, Global is the closest transfer candidate on the pooled panel (Table~\ref{tab:w}). Holding the ER family fixed and replacing development density with each record's own density reduces mean efficiency $W_1$ from 0.374 to 0.153, with a decrease on every one of the eleven records. Relative to Global (0.119) this closes 87\% of the fixed-ER gap; measured from fixed WS the figure is 86\%. The ratio is neither variance explained nor an accuracy: it exceeds 100\% for the school record, where the diagnostic is closer than Global.

\begin{table}[t]\centering
\caption{Efficiency $W_1$ at 50 nodes. \emph{Equal record} gives Facebook ten of eleven votes; \emph{equal source} weights the two source means equally. Target ER uses each test record's density and is diagnostic only.}\label{tab:w}
\small\setlength{\tabcolsep}{3.5pt}
\begin{tabular}{lrrrr}\toprule Target & Global & Fixed ER & Fixed WS & Target ER\\\midrule
Facebook (10)&.117&.365&.356&.156\\School (1)&.138&.463&.455&.121\\Equal record&.119&.374&.365&.153\\Equal source&.128&.414&.406&.139\\\bottomrule\end{tabular}\end{table}

The reason is visible in the candidates' densities. The fixed baselines inherit a density of 0.189 from small, nearly complete development networks, whereas 50-node draws from the larger test records have mean density between 0.03 and 0.14 and Global graphs average 0.061. Across the 77 record--candidate pairs, efficiency distance and density distance are strongly rank-correlated (Spearman $\rho=0.85$; $\rho=0.72$ for the largest-component fraction and $0.09$ for clustering). Much of what the efficiency comparison rewards is having roughly the right number of edges.

Density does not account for every record, and clustering gives a different answer: no candidate is closest at every sample size and source (Figure~\ref{f15}; Appendix~\ref{app:ext}).

\subsection{RQ4: sampling changes the target}\label{sec:cns}
The Copenhagen Networks Study release \citep{cns} is a later structural stress test, not a twelfth row of Table~\ref{tab:w}. Removing eleven self-links from 6,429 rows leaves 6,418 undirected edges among 800 people who appear in the friendship file.

In the full graph, clustering is 0.315 and efficiency 0.364. In 32 uniform draws of 50 endpoints the means are 0.080 and 0.049, while density stays at 0.020 (Figure~\ref{f18}). On average 23.7 of the 50 sampled nodes are isolates. Because nodes of degree zero or one contribute zero local clustering, $\bar C=\Pr(d\ge2)\,E[C_i\mid d_i\ge2]$, so the drop alone does not show weaker closure among nodes that keep two neighbours.

Against these draws (Appendix Table~\ref{tab:cns}), fixed ER is closer than Global on clustering (0.109 versus 0.170). Giving ER the target density reduces its efficiency distance to 0.012, against 0.190 for Global, but its transitivity distance is worse (0.221 versus 0.066). Agreement on reachability is not agreement on triadic closure, and agreement with a sparse induced sample is not agreement with an intact community.

\subsection{Exploratory checks on each question}\label{sec:robust}
Six further analyses, run after the results above, reuse the same graphs and records.

\begin{figure*}[t]\centering
\includegraphics[width=.7\linewidth]{20_sampling_scheme.pdf}
\caption{Mean $W_1$ over the eleven test records when the observed 50-node samples are uniform induced draws or snowball draws (breadth-first from a random seed node, 32 draws per record). Generated graphs and fixed baselines are identical in the two columns; only the observed target changes.}\label{f20}\end{figure*}

\paragraph{RQ4: sampling scheme.} We redrew every test record by snowball sampling: breadth-first search from a uniformly chosen node until 50 nodes are reached, restarting from a new node if a component is exhausted \citep{sampling}. Averaged over records, uniform draws have density 0.083, clustering 0.383 and efficiency 0.197; snowball draws have 0.320, 0.651 and 0.594. The candidates do not change, but their ranking does (Figure~\ref{f20}). Global is closest on efficiency in nine of eleven records under uniform draws and in none under snowball draws, where its mean distance (0.357) is the largest of the seven candidates and fixed ER is closest in eight records (0.131). Iterative, the farthest of the seven candidates on clustering under uniform draws (0.275), becomes closest in nine records (0.077). Neither scheme is the correct one; a fifty-person roster does not say whether it stands for a random set of people or for a connected neighbourhood, and the answer decides the ranking.

\paragraph{RQ3: all six statistics under the density diagnostic.} Section~\ref{sec:density} applied the density-informed ER diagnostic to efficiency. Extending it to every statistic separates what density carries from what it does not. Its mean distance is 0.342 on largest-component fraction (Global 0.340, Iterative 0.339) and 0.196 on degree Gini (best transfer candidate 0.165), so density alone comes close to the best candidate on reachability and degree inequality. It is the worst of all candidates on clustering (0.309; fixed WS 0.122) and transitivity (0.426; Local 0.109), and beats the best transfer candidate on neither in any record. Density explains agreement on reachability and none of the agreement on closure.

\paragraph{RQ2: mixing conditional on the other attributes.} For each nonempty graph and each attribute we permuted that attribute's labels 500 times within strata defined jointly by the other two attributes, which preserves the roster's association between religion and politics, and recorded the observed same-group edge share minus the permutation mean. This conditional excess is positive in all 28 cells for every attribute and procedure. Politics is first in 28 of 28 cells for Local (0.310 against 0.204 for religion), Sequential (0.280 against 0.202) and Iterative (0.385 against 0.110). For Global the order reverses in 28 of 28 cells: religion 0.496, politics 0.251. The procedure effect of Section~\ref{sec:mixing} is therefore not produced by the roster's attribute association, and under this reference the Global ordering is as consistent as the person-level one, in the opposite direction.

\begin{table}[t]\centering
\caption{Joint dyadic model: mean log-odds of a tie for a shared attribute, with 95\% block-bootstrap intervals over the eight schedule blocks. \emph{Gap} is politics minus religion. \emph{Cells} counts the model--setting cells (of 28) in which politics has the largest coefficient.}\label{tab:dyad}
\small\setlength{\tabcolsep}{3pt}
\begin{tabular}{lrrrr}\toprule Procedure & Politics & Religion & Gap [95\% CI] & Cells\\\midrule
Global&2.16&3.79&$-1.63$ $[-1.88,-1.35]$&0\\
Local&2.67&2.24&$0.43$ $[0.29,0.58]$&22\\
Sequential&2.34&2.09&$0.25$ $[0.15,0.35]$&20\\
Iterative&3.51&1.48&$2.03$ $[1.82,2.28]$&28\\\bottomrule\end{tabular}\end{table}

\paragraph{RQ2: a joint dyadic model.} For each of the 878 nonempty graphs we regressed tie presence over the 1{,}225 pairs on shared politics, shared religion, shared gender and the age gap (Table~\ref{tab:dyad}). All three shared-attribute coefficients are positive in all 28 cells for every procedure (gender is smallest, 0.58 to 1.59), and a ten-year age gap multiplies the odds of a tie by 0.20 to 0.40. The politics-minus-religion gap has the same sign in eight of eight blocks for every procedure: negative for Global, positive for the three person-level procedures. The joint model also qualifies the count-based result. Politics has the largest coefficient in 28 of 28 cells only for Iterative; for Local and Sequential it leads in 22 and 20 cells, and their gaps are a fifth to an eighth of Iterative's. The reversal between Global and person-level generation is robust; the margin by which politics leads religion is large only when nominations are revised.

\paragraph{RQ1: tie identity across settings.} A summary statistic can stay fixed while the ties themselves change. We therefore compared edge sets directly, using Jaccard overlap between runs on different schedule blocks. Two runs of the same model, procedure and setting share 0.32 to 0.41 of their ties. A reference run and a run under another setting share 0.99 to 1.01 of that within-setting value for Global, Local and Iterative in all six settings. Sequential retains 0.99 to 1.00 under the three country framings and 0.91, 0.94 and 0.99 under Hindi, Japanese and Portuguese instructions. Changing the framing leaves who is tied to whom as it was, up to run-to-run variation; instruction language moves Sequential ties modestly.

\paragraph{RQ1: what the planned contrasts could detect.} Absence of significant contrasts is informative only if the design could have detected something. Using each contrast's observed block variability, a paired test at the familywise level has 80\% power for a clustering difference of 0.029 for Local, 0.045 for Sequential and 0.051 for Iterative (medians over six contrasts), which is 6\% to 9\% of the reference mean. For Global the detectable density difference is 0.031 against a reference mean of 0.059, so only a change of about half the reference density would have been found; the largest observed Global estimate is 0.010. The person-level nulls rule out large effects. The Global nulls do not.

\section{Discussion and limitations}
A single realism score conflates three choices: who is asked to nominate, what information a baseline may use, and which observed people survive sampling. Our most stable within-design finding is that the nomination procedure reverses the order of demographic mixing, while prompt framing leaves the ties themselves almost untouched. Our strongest external finding is a counterexample to reading short paths as social realism: a baseline with no demographic information recovers most of an apparent efficiency advantage once it is given density.

The limitations are substantial. One fictional roster, a fixed field order and English biographies prevent attribution to demographic preferences or to national culture; the stratified permutation removes the roster's attribute association as an explanation within this roster but cannot say what another roster would give. Procedures were not compute-matched: the median ratio of requests to Global was about 50 for Local and Sequential and 347 for Iterative. The four configurations are a fixed panel, not a sample of models. Conditioning on valid output can favour procedures that fail on harder inputs; missingness bounds and complete-block counts describe this problem without removing it. The ten Facebook records can share people, the school survey misses some nominations, and the Copenhagen file omits people without observed ties. Test records were held out of baseline fitting, not necessarily out of language-model training data. The density split, the rank correlation, the Copenhagen analysis and the six analyses of Section~\ref{sec:robust} were not pre-specified. Block-bootstrap intervals rest on eight blocks and are likely too narrow. Similarity of graph statistics across unaligned people cannot validate prediction of individual links.

The remaining gaps need new generation: vary roster and field order independently, add a shuffled-attribute control in the spirit of \citet{chang}, and match computation across procedures. A confirmatory external comparison should declare a density-matched candidate and fix the sampling scheme in advance.

\section{Conclusion}
Prompt framing left the generated ties largely unchanged; the nomination procedure reversed which attribute organised them; a density-informed random graph recovered most of an apparent reachability advantage and none of the closure; and the sampling scheme decided which generator was closest. A network match becomes interpretable when these four checks are reported with it.

\section*{Impact Statement}
Synthetic networks can expose weaknesses in evaluation without exposing identifiable respondents. Conversely, treating a generated mixing pattern as a fact about religious, political or language communities could reinforce stereotypes. We report aggregate structural comparisons on fictional personas and publicly released, de-identified network data, and we do not recommend using these results to infer any person's attributes or relationships.

\bibliography{references}
\bibliographystyle{icml2026}

\clearpage\onecolumn\appendix
\section{Primary estimates}
\begin{table}[ht]\centering
\caption{All 24 planned contrasts: difference from US-English in density (Global) or mean local clustering (other procedures). Intervals are Bonferroni simultaneous 95\% intervals; Holm $p$-values are adjusted over the same family of 24. Dashes mark the contrast with fewer than six complete blocks.}\label{contrasts}
\small\begin{tabular}{llrrrrr}\toprule Procedure & Framing / instruction language & Blocks & Difference & Simultaneous 95\% interval & Raw $p$ & Holm $p$\\\midrule
Global & India / English & 8 & $-0.0030$ & $[-0.0300,\ +0.0239]$ & 0.612 & 1 \\
Global & Japan / English & 8 & $-0.0022$ & $[-0.0251,\ +0.0207]$ & 0.663 & 1 \\
Global & Brazil / English & 8 & $-0.0012$ & $[-0.0273,\ +0.0249]$ & 0.833 & 1 \\
Global & US / Hindi & 8 & $+0.0097$ & $[-0.0078,\ +0.0272]$ & 0.0342 & 0.616 \\
Global & US / Japanese & 5 & $+0.0036$ & -- & -- & -- \\
Global & US / Portuguese & 6 & $+0.0066$ & $[-0.0499,\ +0.0630]$ & 0.527 & 1 \\
\addlinespace
Local & India / English & 8 & $+0.0061$ & $[-0.0195,\ +0.0316]$ & 0.296 & 1 \\
Local & Japan / English & 8 & $-0.0132$ & $[-0.0326,\ +0.0062]$ & 0.0144 & 0.288 \\
Local & Brazil / English & 8 & $+0.0022$ & $[-0.0291,\ +0.0335]$ & 0.749 & 1 \\
Local & US / Hindi & 8 & $-0.0171$ & $[-0.0398,\ +0.0056]$ & 0.0090 & 0.207 \\
Local & US / Japanese & 8 & $-0.0180$ & $[-0.0438,\ +0.0078]$ & 0.0128 & 0.268 \\
Local & US / Portuguese & 8 & $-0.0053$ & $[-0.0292,\ +0.0186]$ & 0.330 & 1 \\
\addlinespace
Sequential & India / English & 8 & $-0.0028$ & $[-0.0561,\ +0.0504]$ & 0.807 & 1 \\
Sequential & Japan / English & 8 & $0.0000$ & $[-0.0331,\ +0.0330]$ & 0.996 & 1 \\
Sequential & Brazil / English & 8 & $-0.0045$ & $[-0.0360,\ +0.0270]$ & 0.522 & 1 \\
Sequential & US / Hindi & 8 & $-0.0001$ & $[-0.0420,\ +0.0419]$ & 0.996 & 1 \\
Sequential & US / Japanese & 8 & $-0.0171$ & $[-0.0443,\ +0.0102]$ & 0.0207 & 0.394 \\
Sequential & US / Portuguese & 8 & $-0.0129$ & $[-0.0654,\ +0.0397]$ & 0.282 & 1 \\
\addlinespace
Iterative & India / English & 8 & $-0.0069$ & $[-0.0604,\ +0.0465]$ & 0.557 & 1 \\
Iterative & Japan / English & 8 & $-0.0016$ & $[-0.0407,\ +0.0374]$ & 0.849 & 1 \\
Iterative & Brazil / English & 8 & $+0.0017$ & $[-0.0305,\ +0.0339]$ & 0.808 & 1 \\
Iterative & US / Hindi & 8 & $-0.0453$ & $[-0.0850,\ -0.0056]$ & 0.0010 & 0.0238 \\
Iterative & US / Japanese & 7 & $-0.0532$ & $[-0.1286,\ +0.0221]$ & 0.0107 & 0.235 \\
Iterative & US / Portuguese & 8 & $-0.0203$ & $[-0.0664,\ +0.0258]$ & 0.0749 & 1 \\
\bottomrule\end{tabular}\end{table}

\begin{table}[ht]\centering
\caption{Copenhagen: $W_1$ distance to 50-node induced draws for all six statistics. Each column has its own units and columns are not combined. The final row is target-informed and is not a transfer candidate.}\label{tab:cns}
\small\begin{tabular}{lrrrrrr}\toprule Candidate & Density & Clustering & Efficiency & Transitivity & Degree Gini & Largest component\\\midrule
Global&.042&.170&.190&.066&.488&.535\\
Local&.149&.442&.460&.192&.339&.721\\
Sequential&.151&.436&.470&.190&.307&.722\\
Iterative&.153&.572&.362&.426&.303&.613\\
Fixed ER&.169&.109&.521&.097&.458&.722\\
Fixed BA&.163&.215&.520&.097&.343&.722\\
Fixed WS&.184&.319&.513&.163&.548&.722\\\midrule
Target ER (diagnostic)&.002&.074&.012&.221&.104&.044\\\bottomrule\end{tabular}\end{table}

\section{Evidence boundary and reproducibility}
All estimates use the primary collection only. Nine separately labelled supplementary attempts yielded six graphs and three failures; none enters a denominator, model fit, rank or figure. The planned contrast table was fixed before the held-out comparison; the mixing, density-split, rank-correlation, Copenhagen and Section~\ref{sec:robust} analyses were added afterwards and are exploratory. An anonymized snapshot of the code, the released numerical tables and the executed analysis notebooks accompanies this submission. From that snapshot, \texttt{python check\_efficiency.py .} reads the eleven released record-level rows, recomputes both gap ratios and the weighting rows of Table~\ref{tab:w}, and recomputes the rank correlations of Section~\ref{sec:density}; it makes no network or model calls. \texttt{python scripts/run\_robustness\_checks.py --output <new directory>} regenerates every number and the figure of Section~\ref{sec:robust} from the primary graphs and the public records in about one minute, with fixed seeds (32 snowball draws per record, eight banks of 56 graphs for the diagnostic, 500 permutations per graph); it asserts that its efficiency column reproduces the released table and refuses to write into an existing directory. \texttt{python scripts/run\_tie\_level\_checks.py --output <new directory>} does the same for the dyadic model (ridge-penalised logistic regression, $C=10$, 10{,}000 block-bootstrap draws), the tie-overlap comparison and the detectable differences, in under ten seconds. The 25- and 40-node clustering margins are in the released matched-size table. Public observations are obtained from their original sources and are not redistributed as person-level records.

\section{External comparison: record-level detail}\label{app:ext}
Density does not account for every record (Figure~\ref{f15}). In a post hoc split at full-record density 0.06, the five sparser records favour the diagnostic on average (0.110 versus 0.148 for Global) and the six denser records favour Global (0.095 versus 0.189); FB02 and FB05 run against their groups by small margins. Eleven overlapping records do not support a fitted trend.

Clustering gives a different answer. At 50 nodes WS is closest on average for Facebook (0.126) and BA for the school (0.042). At 40 nodes on Facebook the WS--BA margin is 0.004, and at 25 nodes on the school record ER is closer than Global by 0.017. Global's clustering distance is 0.096 on the sparsest record and 0.417 on one of the densest. No candidate is closest throughout (Appendix Figures~\ref{f08} and~\ref{f12}). Mean observed clustering exceeds a degree-preserving rewired control in each of the eleven records (Appendix Figure~\ref{f11}), so fixing degrees does not reproduce closure; finite edge swaps do not guarantee a uniform null sample, and we make no test of that difference.

\section{Additional figures}
Figures~\ref{f02}--\ref{f05} document planned inference, failures and resources; Figures~\ref{f06}--\ref{f13} cover development context, transfer distances, conditional controls and size sensitivity; Figure~\ref{f17} gives group-pair mixing; Figure~\ref{f19} gives the Copenhagen comparison.

\begin{figure}[htbp]\centering
\includegraphics[width=.9\linewidth,height=.30\textheight,keepaspectratio]{02_primary_robustness.pdf}
\caption{Iterative US-Hindi versus US-English. (a) The eight paired block differences with the simultaneous interval. (b) Raw (crosses) and Holm-adjusted (circles) $p$ after omitting each block; the refits overlap and are not replications. (c) The $p$-values an eight-block exact sign-flip test can return, against the first Holm threshold.}\label{f02}\end{figure}
\begin{figure}[htbp]\centering
\includegraphics[width=.8\linewidth,height=.42\textheight,keepaspectratio]{03_model_profiles.pdf}
\caption{Model-specific mean paired differences from US-English on common schedule blocks, with the four-model mean and the number of models sharing its sign. Descriptive only; no intervals or tests.}\label{f03}\end{figure}
\begin{figure}[htbp]\centering
\includegraphics[width=.9\linewidth,height=.34\textheight,keepaspectratio]{04_coverage_missingness.pdf}
\caption{Every primary attempt as one square (valid nonempty, valid empty, failed), and identification bounds for the three contrasts that lost blocks. Bounds are not confidence intervals; missing outcomes are never set to zero.}\label{f04}\end{figure}
\begin{figure}[htbp]\centering
\includegraphics[width=.9\linewidth,height=.26\textheight,keepaspectratio]{05_resource_use.pdf}
\caption{Requests, estimated cost and provider time per attempt on log scales. Procedures were not compute-matched.}\label{f05}\end{figure}
\begin{figure}[htbp]\centering
\includegraphics[width=.9\linewidth,height=.36\textheight,keepaspectratio]{06_gnm_context.pdf}
\caption{Observed mean clustering minus the mean over eight random graphs with the same nodes and edge count, for every valid graph. A conditional structural reference, not a transfer prediction.}\label{f06}\end{figure}
\begin{figure}[htbp]\centering
\includegraphics[width=.9\linewidth,height=.30\textheight,keepaspectratio]{07_legacy_comparison.pdf}
\caption{Valid generated graphs beside the eight archived development networks. These networks fixed the baselines and are not test outcomes; their sizes and tie definitions differ.}\label{f07}\end{figure}
\begin{figure}[htbp]\centering
\includegraphics[width=.9\linewidth,height=.40\textheight,keepaspectratio]{08_heldout_distances.pdf}
\caption{$W_1$ distance to 50-node draws from the eleven held-out records, one panel per statistic. Ticks are the ten Facebook records, circles their mean, diamonds the school record. Panels have different units and are not combined.}\label{f08}\end{figure}
\begin{figure}[htbp]\centering
\includegraphics[width=.9\linewidth,height=.36\textheight,keepaspectratio]{09_heldout_distributions.pdf}
\caption{Distributions behind the distances: density, clustering and efficiency of observed 50-node draws and of every candidate, with medians, interquartile ranges and 5th--95th percentile spans.}\label{f09}\end{figure}
\begin{figure}[htbp]\centering
\includegraphics[width=.9\linewidth,height=.32\textheight,keepaspectratio]{10_sampling_effect.pdf}
\caption{Mean of 32 induced 50-node draws divided by the full-record value, against the share of the record in a draw. Lines span the draws and are not confidence intervals.}\label{f10}\end{figure}
\begin{figure}[htbp]\centering
\includegraphics[width=.9\linewidth,height=.38\textheight,keepaspectratio]{11_conditional_controls.pdf}
\caption{Observed draw minus controls rebuilt from that draw: same edge count (squares) and same degree sequence (circles). Controls use target information and are diagnostics.}\label{f11}\end{figure}
\begin{figure}[htbp]\centering
\includegraphics[width=.9\linewidth,height=.40\textheight,keepaspectratio]{12_cohort_sensitivity.pdf}
\caption{Clustering $W_1$ by record at 50 nodes, ordered by full-record density, with the gap between the closest and second-closest candidate. Lines join separate records only to show changes of order.}\label{f12}\end{figure}
\begin{figure}[htbp]\centering
\includegraphics[width=.9\linewidth,height=.34\textheight,keepaspectratio]{13_appendix_lcc.pdf}
\caption{Largest-component fraction. Nearly every ER, BA and WS graph is connected, so their distances on this statistic are almost identical.}\label{f13}\end{figure}
\begin{figure}[htbp]\centering
\includegraphics[width=.85\linewidth,height=.48\textheight,keepaspectratio]{17_rq2_mixing_matrices.pdf}
\caption{Group-pair share of tie ends divided by the roster-chance share (1 is chance), by attribute and procedure. Colour compares pairs within one attribute and procedure only. In Iterative, the left--right entries shown as 0.0 are a rounded ratio of 0.044, not an absence of ties.}\label{f17}\end{figure}
\begin{figure}[htbp]\centering
\includegraphics[width=.9\linewidth,height=.42\textheight,keepaspectratio]{19_copenhagen_transfer.pdf}
\caption{Copenhagen at 50 nodes: $W_1$ for five statistics; Table~\ref{tab:cns} also reports the largest-component fraction. Ticks are individual models or simulation banks. The shaded row is the target-informed diagnostic.}\label{f19}\end{figure}
\end{document}
